% Prostudováno: ucebnice/02_Funkce_a_rovnice/Zaklady_funkci/Funkce.tex, části D/H, monotonie, sudost a lichost;
% ucebnice/02_Funkce_a_rovnice/Linearni_funkce_a_rovnice/Lineární funkce.tex, předpis a graf lineární funkce.
% Realisticky.cz: externi_zdroje/realisticky/02_Funkce_a_rovnice/01_Linearni_funkce/020105_Graf_funkce_I/020105_Lekce.pdf
% a 020113_Rostouci_a_klesajici_funkce_I/020113_Lekce.pdf; metodika přechodu mezi osami a zdůvodnění růstu. Úlohy vlastní.
\documentclass[11pt,a4paper]{article}
% Spolecny styl pro ucitelske pripravy MAT (Work/MAT)
% Verze 1.1, 2026-10-23
\NeedsTeXFormat{LaTeX2e}
\ProvidesPackage{priprava}[2026/10/23 v1.1 Sablona ucitelskych priprav MAT]

% Zakladni typografie a cestina
\RequirePackage[T1]{fontenc}
\RequirePackage[utf8]{inputenc}
\RequirePackage[czech]{babel}
\RequirePackage{lmodern}
\RequirePackage{microtype}
\RequirePackage{amsmath,amssymb}
\RequirePackage{xcolor}
\RequirePackage{enumitem}
\RequirePackage[most]{tcolorbox}
\RequirePackage{tabularx,booktabs,array}
\RequirePackage{fancyhdr}
\RequirePackage{needspace}
\RequirePackage{geometry}
\RequirePackage[hidelinks]{hyperref}

\geometry{a4paper,top=17mm,bottom=18mm,left=23mm,right=23mm}
\definecolor{notegray}{gray}{0.39}
\setlength{\parindent}{0pt}
\setlength{\parskip}{2.5pt plus 1pt minus 0.5pt}
\setlist[itemize]{topsep=2pt,itemsep=1pt,parsep=0pt,leftmargin=1.45em}
\setlist[enumerate]{topsep=2pt,itemsep=2pt,parsep=0pt,leftmargin=1.65em}
\setcounter{secnumdepth}{0}
\renewcommand{\arraystretch}{1.12}

% Zahlavi: vlevo rocnik | globalni tema, vpravo cislo: nazev hodiny
% Zapati zustava zcela prazdne.
\pagestyle{fancy}
\fancyhf{}
\renewcommand{\headrulewidth}{0.35pt}
\renewcommand{\footrulewidth}{0pt}
\setlength{\headheight}{14pt}
\newcommand{\nastavPripravu}[4]{%
  \fancyhead[L]{\small #1\enspace\textbar\enspace #2}%
  \fancyhead[R]{\small #3: #4}%
}

% Minimalisticke cernobile boxy. Kazda definice zvlast.
\tcbset{commonlesson/.style={enhanced,breakable=false,colback=white,colframe=black,
  colbacktitle=white,coltitle=black,boxrule=0.45pt,
  arc=0.8mm,outer arc=0.8mm,left=3mm,right=3mm,top=1.6mm,bottom=1.6mm,
  toptitle=1mm,bottomtitle=0.7mm,fonttitle=\bfseries,
  attach boxed title to top left={xshift=2.5mm,yshift=-1.15mm},
  boxed title style={colback=white,colframe=white,boxrule=0pt,left=1mm,right=1mm}}}
\newtcolorbox{zadaniBox}[1][]{commonlesson,title={Zadání},#1}
\newtcolorbox{prikladBox}[1][]{commonlesson,title={Příklad},#1}
\newtcolorbox{definiceBox}[1][]{commonlesson,title={Definice},#1}

% Spolecny box pro zadani a jeho reseni, bez mezery mezi castmi.
% Pouziti:
%   \begin{prikladReseniBox}
%     Zadani ulohy.
%     \tcblower
%     \begin{align*} ... \end{align*}
%   \end{prikladReseniBox}
% Spodni cast ma jemny sedy podklad a delici linku; lze v ni sazet
% vice radku matematiky. Obe casti maji jen jeden spolecny ramecek.
\tcbset{solutionlesson/.style={
  commonlesson,
  enhanced jigsaw,
  breakable,
  colbacklower=black!3,
  segmentation style={solid,draw=black!25,line width=0.35pt},
  fontlower=\small\color{notegray},
  before lower={\noindent\textbf{Řešení:}\enspace},
  leftlower=3mm,rightlower=3mm
}}
\newtcolorbox{prikladReseniBox}[1][]{solutionlesson,title={Příklad},#1}
\newtcolorbox{zadaniReseniBox}[1][]{solutionlesson,title={Zadání},#1}

% Pomucky pro autora pripravy: sedive poznamky; \reseni pro starsi pripravy
\newcommand{\poznamka}[1]{\par\begingroup\small\color{notegray}\itshape
  \textnormal{\textbf{Poznámka pro učitele.}}\enspace #1\par\endgroup}
\newcommand{\reseni}[1]{\par\begingroup\small\color{notegray}
  \textbf{Očekávané řešení:}\enspace #1\par\endgroup}
\newcommand{\blok}[1]{\needspace{5\baselineskip}\subsection*{#1}}
\newcommand{\vyrok}[1]{\uv{#1}}

\usepackage{tikz}
\nastavPripravu{1. ročník}{Funkce}{71}{Procvičování grafů a vlastností}
\begin{document}
\blok{Samostatné procvičování s cílenou podporou}
Hodina navazuje na grafy lineárních a konstantních funkcí i na nový graf absolutní hodnoty. Žáci nejprve sami převedou předpis na tabulku a graf, potom z grafu zpět vyčtou vlastnosti. Učitel před rozdáním úloh ukáže měřítko os a připomene, že definiční obor určuje, kterou část přímky skutečně kreslíme.
\begin{definiceBox}
\textbf{Graf funkce} $f$ je množina bodů $(x;f(x))$ pro všechna $x\in D(f)$. Bod patří do grafu právě tehdy, když jeho první souřadnice patří do $D(f)$ a druhá se rovná $f(x)$.
\end{definiceBox}
\begin{prikladReseniBox}
Na oboru $\langle-2;3)$ je dána funkce $f(x)=x+1$. Sestrojte graf, určete $D(f)$, $H(f)$, průsečíky s osami a monotonii.
\tcblower
Tabulka pro $x=-2,-1,0,2$ dává $f(x)=-1,0,1,3$. Kreslíme část přímky od plného bodu $(-2;-1)$ k prázdnému bodu $(3;4)$. Proto $D(f)=\langle-2;3)$ a $H(f)=\langle-1;4)$. Průsečíky jsou $(-1;0)$ a $(0;1)$. Pro $x_1<x_2$ je $f(x_2)-f(x_1)=x_2-x_1>0$, tedy funkce roste na celém oboru.
\end{prikladReseniBox}
\begin{center}
\begin{tikzpicture}[x=.68cm,y=.58cm,>=stealth]
\draw[step=1,gray!25,very thin] (-3,-2) grid (4,5);
\draw[->] (-3.3,0)--(4.3,0) node[right] {$x$};\draw[->] (0,-2.3)--(0,5.3) node[above] {$y$};
\foreach \x in {-2,-1,1,2,3}{\draw (\x,.08)--(\x,-.08) node[below] {\small $\x$};}
\foreach \y in {-1,1,2,4}{\draw (.08,\y)--(-.08,\y) node[left] {\small $\y$};}
\draw[very thick] (-2,-1)--(3,4);
\fill (-2,-1) circle (2.4pt);\draw[fill=white,thick] (3,4) circle (2.7pt);
\node[below left] at (-2,-1) {\small plný};\node[above right] at (3,4) {\small prázdný};
\end{tikzpicture}
\end{center}
\poznamka{V hlavním procvičení neprozrazuj celý postup. Vzor slouží jen k připomenutí otevřeného konce. Žák musí vždy říci, zda určuje vlastnost předpisu na $\mathbb R$, nebo skutečně zadané funkce na omezeném oboru.}
\blok{Jádro samostatné práce}
\begin{zadaniReseniBox}
\begin{enumerate}[label=\alph*)]
\item Pro $g(x)=-2x+2$ na $\mathbb R$ sestavte tabulku pro $x=-1,0,1,2$, sestrojte přímku a určete průsečíky s osami i směr změny.
\item Pro $h(x)=3$ na $\langle-1;2\rangle$ sestrojte graf, určete obor hodnot a rozhodněte, zda je funkce rostoucí, klesající, neklesající či nerostoucí.
\item Z grafu určeného úsečkami přes $A(-2;1),B(0;-1),C(2;1)$ určete $D,H$, nulové hodnoty, intervaly růstu a poklesu. Ověřte jedním výpočtem hodnotu v $x=-1$.
\end{enumerate}
\tcblower
\begin{enumerate}[label=\alph*)]
\item Hodnoty $4,2,0,-2$; přímka prochází $(0;2)$ a $(1;0)$. Při zvýšení $x$ o 1 klesne hodnota o 2, tedy je klesající; průsečíky jsou $(0;2)$ a $(1;0)$.
\item Graf je vodorovná uzavřená úsečka od $(-1;3)$ do $(2;3)$; $H(h)=\{3\}$. Není přísně rostoucí ani klesající; je současně neklesající i nerostoucí, protože pro libovolná $x_1<x_2$ vyjde $h(x_1)=h(x_2)$.
\item $D=\langle-2;2\rangle$, $H=\langle-1;1\rangle$. Nuly jsou $x=-1$ a $x=1$. Funkce klesá na $\langle-2;0\rangle$ a roste na $\langle0;2\rangle$. První úsečka má sklon $-1$ a předpis $y=-x-1$, tedy $f(-1)=0$.
\end{enumerate}
\end{zadaniReseniBox}
\blok{Kontrola porozumění a argumentace}
\begin{zadaniReseniBox}
\begin{enumerate}[label=\alph*)]
\item Mohou mít různé lineární funkce stejný průsečík s osou $y$? Uveďte dvě a vysvětlete parametr.
\item Žák tvrdí: „Z tabulky $(-1;1),(0;0),(1;1)$ plyne, že funkce je $x^2$.“ Posuďte závěr.
\item Rozhodněte, zda $r(x)=2x$ na $\mathbb R$ je sudá, lichá, nebo ani jedna; zdůvodněte předpisem i souměrností grafu.
\end{enumerate}
\tcblower
\begin{enumerate}[label=\alph*)]
\item Ano, např. $y=x+2$ a $y=-2x+2$ mají společnou hodnotu $b=2$ při $x=0$, ale různé směrnice.
\item Neplyne. Stejné tři hodnoty má například $x^2+x(x-1)(x+1)$; přidaný součin je v $-1,0,1$ nulový, jinde se předpisy liší. Bez další podmínky tabulka obecný předpis neurčuje.
\item $r(-x)=-2x=-r(x)$ pro všechna $x\in\mathbb R$, takže je lichá. Graf $y=2x$ je souměrný podle počátku.
\end{enumerate}
\end{zadaniReseniBox}
\blok{Volitelné rozšíření a rezerva}
\begin{zadaniReseniBox}
Pro $p(x)=|x-1|$ na $\langle-1;3)$ sestrojte graf, určete $D,H$, minimum a monotónní části. Jak se změní obor hodnot po nahrazení pravého konce intervalem uzavřeným?
\tcblower
Vrchol $(1;0)$ do oboru patří. Vlevo je $p(-1)=2$, vpravo se hodnota blíží $2$ a koncový bod chybí; hodnota $2$ však již nastává vlevo. Tedy $H=\langle0;2\rangle$. Klesá na $\langle-1;1\rangle$, roste na $\langle1;3)$ a minimum je $0$. Uzavření pravého konce obor hodnot nezmění.
\end{zadaniReseniBox}
\begin{zadaniReseniBox}
Najděte lineární funkci s nulovou hodnotou při $x=3$ a s hodnotou $2$ při $x=1$.
\tcblower
Směrnice je $(0-2)/(3-1)=-1$. Z $2=-1\cdot1+b$ plyne $b=3$, tedy $f(x)=-x+3$. Kontrola: $f(3)=0$ a $f(1)=2$.
\end{zadaniReseniBox}
\blok{Domácí práce}
\begin{zadaniReseniBox}
\begin{enumerate}[label=\alph*)]
\item Sestrojte $f(x)=\frac12x-1$ na $\mathbb R$ pomocí tří bodů a určete průsečíky s osami.
\item Pro $g(x)=2-x$ na $\langle0;4\rangle$ určete $D,H$, krajní body, nulovou hodnotu a monotonii.
\item Vysvětlete, proč funkce $h(x)=|x|$ na $\mathbb R$ není rostoucí na celém oboru, přestože pro $x>0$ roste.
\end{enumerate}
\tcblower
\begin{enumerate}[label=\alph*)]
\item Např. pro $x=0,2,4$ jsou hodnoty $-1,0,1$; průsečíky $(0;-1)$ a $(2;0)$.
\item Krajní body $(0;2),(4;-2)$; $D=\langle0;4\rangle$, $H=\langle-2;2\rangle$, nula v $x=2$. Směrnice $-1$ značí pokles.
\item Pro $-2<-1$ dostaneme $h(-2)=2>1=h(-1)$, což odporuje požadavku růstu na celé $\mathbb R$. Na $\langle0;\infty)$ roste.
\end{enumerate}
\end{zadaniReseniBox}
\poznamka{Jádro 45 min: organizace 3, rychlé připomenutí a vzor 7, samostatná práce 22, společná kontrola a oprava argumentace 13. Rezerva je mimo součet. Domácí práce přibližně 15 minut. U slabší skupiny dej úlohu c) z jádra až po první kontrole.}
\end{document}
